% IEEE-format project report: "A Paired Regression Gate for LLM Features".
% The text mirrors docs/report/01-abstract.md ... 08-conclusion.md; keep the two in step when editing.
% Tables in tables/ are generated by scripts/make_report_tables.py from experiments/*/expected/*.json:
% never edit them by hand. Build from this directory: pdflatex main && bibtex main && pdflatex main && pdflatex main
\documentclass[conference]{IEEEtran}
\usepackage[T1]{fontenc}
\usepackage[utf8]{inputenc}
\usepackage{amsmath,amssymb}
\usepackage{array}
\usepackage{url}
\usepackage{cite}
\emergencystretch=1.5em

\begin{document}

\title{A Paired Regression Gate for LLM Features:\\ Execution-Grounded Scoring and Recorded Replies\\ in Continuous Integration}

% PLACEHOLDER: replace with the author block before submission.
\author{\IEEEauthorblockN{Author name (placeholder)}
\IEEEauthorblockA{Course and institution (placeholder)}}

\maketitle

\begin{abstract}

A feature built on a large language model (LLM) changes behaviour whenever its prompt, model, adapter, weight quantization or serving engine changes, and such a change can degrade the feature without any code change that a conventional test would notice. This report presents an evaluation harness that decides, by a written and tested rule, whether a candidate configuration of an LLM feature may replace the current baseline, and that enforces the decision in continuous integration (CI). The feature under test is an LLM bug-fix assistant. Its replies are scored by strict schema validation, an abstract-syntax-tree (AST) based static scan, and execution of shown and hidden tests in a sandboxed child process whose results are compared outside the child. A second feature, a patch-correctness judge, is scored against labels obtained by execution. The gate pairs the two arms on item identifiers and applies a non-inferiority rule with an explicit tie policy, hard validity gates, a new-violation gate for the static scan and an exact McNemar guard; every report carries a paired cluster-bootstrap interval and the minimum detectable effect of its suite. Model replies are recorded once as request-addressed ``cassettes'', so CI replays them without a GPU, network access or secrets and reproduces every report byte for byte, and a deliberately regressed candidate serves as a self-test of the gate. Applied to low-rank (LoRA) adapters on a 27-billion-parameter open-weight model, the gate rejected an identity adapter, which cannot change any output in exact arithmetic, because one reply changed on each bug-fix split at serving time. A trained judge adapter improved balanced accuracy on the development split by +0.154 (4 gains, 0 losses, McNemar p = 0.125) but could not show the gain on a held-out split on which the base model already scored 1.000. Six of the seven items on which the arms disagreed were failures of format or truncation, not wrong answers. At suite sizes of tens of items, a regression gate must therefore be calibrated against a measured no-op noise floor and checked for ceiling effects before its verdicts carry information.

\end{abstract}

\begin{IEEEkeywords}
large language models, evaluation, regression testing, continuous integration, statistical hypothesis testing, low-rank adaptation, quantization
\end{IEEEkeywords}

\section{Introduction}

\subsection{Problem}

Software teams increasingly ship features whose behaviour is produced by an LLM: a prompt template, a model, optional adapters and a serving configuration together determine the output. Each of these is changed routinely. A prompt is reworded, a model is replaced by a newer release, a parameter-efficient adapter is trained~\cite{hu2022lora}, weights are quantized to fit cheaper hardware~\cite{frantar2023gptq,ma2024bitnet}, or the serving engine is replaced~\cite{kwon2023vllm,llamacpp}. None of these changes touches the code that unit tests exercise, yet any of them can make the feature worse. The question that a regression test must answer for such a change is therefore not whether the code is correct, but whether the change makes the feature worse on inputs it was not tuned on, and how confident that answer can be.

Benchmarks for code generation answer a related but different question: how good a model is, on average, on a public task set~\cite{chen2021codex,austin2021mbpp,jimenez2024swebench}. A regression gate instead compares two configurations of one product feature, decides whether the second may replace the first, and must make that decision automatically, reproducibly and cheaply on every change, typically in CI, where calling a 27-billion-parameter model is neither free nor deterministic.

\subsection{Motivation}

Three observations motivate a dedicated harness. First, the comparison is between two arms on the same items, so a paired analysis is available and much stronger than comparing two aggregate scores~\cite{dietterich1998,miller2024}. Second, the suites a product team can afford to label with executable tests are small, tens to hundreds of items, and on such suites statistical power is low~\cite{card2020power}; a gate that does not state how large a change it could detect invites over-interpretation of both its passes and its failures. Third, the decision rule must be explicit about ties. A rule of the form ``candidate >= baseline'' promotes every change that does nothing, and a rule that demands significance rejects nearly every candidate on a small suite. Non-inferiority testing, as used in clinical trials~\cite{walker2011}, makes the margin and the direction of the burden of proof explicit, and the same discipline applies here.

LLM-specific failure modes add requirements. A reply can be malformed, truncated at the token limit, or correct only on the test the model was shown. Generated code can call \texttt{eval} or spawn processes, and a judge model can approve every patch while sounding confident~\cite{zheng2023judge}. Serving is nondeterministic even at temperature zero~\cite{yuan2025nondeterminism,he2025nondeterminism}, so two recordings of the same configuration can differ.

\subsection{Contributions}

This report makes the following contributions, each implemented and tested in the accompanying repository:

\begin{enumerate}
\item \textbf{A paired non-inferiority gate for LLM features} (rule \texttt{noninferior/\allowbreak{}1}, \texttt{src/\allowbreak{}llm\_\allowbreak{}eval\_\allowbreak{}harness/\allowbreak{}gate.\allowbreak{}py}) with an explicit margin and tie policy, hard gates on schema and thinking validity, a gate on newly introduced static-analysis violations, an exact McNemar guard~\cite{mcnemar1947}, and an optional significant-gain requirement. Every report states its own detection limit: the minimum detectable effect at 80\% power and the fewest losses McNemar's test would call significant.
\item \textbf{Execution-grounded scoring} (\texttt{verdicts.\allowbreak{}py}, \texttt{sandbox.\allowbreak{}py}, \texttt{sast.\allowbreak{}py}) in which tests are stored as data, the child process that runs model-written code returns only values, and the parent compares them, so that an exit status of 0, an early exit or a skipped test is never a pass. Shown and hidden tests separate genuine fixes from patches overfitted to the shown test.
\item \textbf{An execution-labelled judge evaluation} in which every task contributes one correct and one plausible wrong patch, so that a judge that ignores the patch scores exactly 0.5 balanced accuracy.
\item \textbf{Record and replay of model replies} (\texttt{cassettes.\allowbreak{}py}, \texttt{clients.\allowbreak{}py}): each reply is stored under the SHA-256 of the exact request that produced it, replay refuses a recording that does not answer today's request, and every report is byte-identical across replays. CI therefore needs no GPU, network or secrets, and it gates a deliberately regressed candidate in every run as a self-test of the gate.
\item \textbf{A hash-chained promotion ledger} (\texttt{ledger.\allowbreak{}py}) that records promotions and rollbacks, re-fingerprints the current baseline in CI, and counts held-out evaluations against a budget.
\item \textbf{An empirical study on a 27-billion-parameter open-weight model} with LoRA adapters (Experiment A), which shows that an identity adapter is rejected by a margin-0 gate because of serving nondeterminism, and that a held-out split at ceiling cannot confirm a development-split gain. A quantization study (Experiment C) is designed with an engine-control arm, and its tables are generated automatically once it is recorded.
\end{enumerate}

\subsection{Organization}

Section II reviews related work. Section III describes the system. Section IV states the experimental methodology and the hypotheses before Section V reports the results. Section VI discusses the results and the threats to their validity, and Section VII concludes.

\section{Related Work}

\subsection{Execution-based evaluation of code generation}

HumanEval~\cite{chen2021codex} established functional correctness, measured by executing unit tests, as the standard for evaluating generated code, and MBPP~\cite{austin2021mbpp} provided a larger set of entry-level Python problems with test cases. HumanEvalPack~\cite{muennighoff2024octopack} extended HumanEval to code repair, explanation and synthesis in six languages; its repair task, HumanEvalFix, gives the model a function with an inserted bug and is the task family closest to the bug-fix feature evaluated here. SWE-bench~\cite{jimenez2024swebench} moved evaluation to real repositories, with issues and pull requests from popular Python projects. EvalPlus~\cite{liu2023evalplus} showed that the original HumanEval tests miss a substantial share of wrong programs, and that adding generated tests lowers reported pass rates. The present harness adopts the same lesson at small scale: every task has hidden tests beyond the one shown to the model, and a patch that passes the shown test but fails a hidden one is counted separately as an overfit. Unlike these benchmarks, the harness is not designed to rank models in general; it compares two configurations of one feature and decides whether one may replace the other.

\subsection{Contamination and the reuse of held-out data}

Public benchmarks may be part of a model's training data. Sainz et al.~\cite{sainz2023contamination} argue that contamination must be measured per benchmark, and LiveCodeBench~\cite{jain2025livecodebench} addresses it by evaluating only on problems published after a model's training cut-off. A second, subtler form of leakage arises from the evaluator's own decisions: each time a held-out set is used to choose between candidates, information about it flows into the choice. Dwork et al.~\cite{dwork2015reusable} formalise this adaptive reuse of a holdout and propose mechanisms that preserve validity, and Recht et al.~\cite{recht2019imagenet}, who rebuilt the CIFAR-10 and ImageNet test sets, found accuracy drops of 3\% to 15\% and 11\% to 14\% on the new sets, which they attribute to slightly harder images rather than to adaptive overfitting; a reused test set is thus not necessarily spent, but whether it is can only be established with fresh data. The harness splits tasks by a hash of a stable key rather than of their text, keeps near-duplicate ``twin'' tasks on one side of the split, and its ledger counts the distinct candidates evaluated on the held-out split and warns once a budget is spent.

\subsection{LLMs as judges}

Zheng et al.~\cite{zheng2023judge} found that strong LLM judges agree with human preferences at rates comparable to inter-human agreement, while documenting position, verbosity and self-enhancement biases. For code, a judge's opinion can be checked against execution. The judge feature of this harness is evaluated only against labels obtained by running tests, on a design in which every task contributes one correct and one wrong patch, so that balanced accuracy separates a judge that reads the patch from one that does not.

\subsection{Statistical comparison of models}

McNemar's test~\cite{mcnemar1947} compares two classifiers evaluated on the same items using only the items on which they disagree; Dietterich~\cite{dietterich1998} found it to have low type I error among tests for comparing learning algorithms. Paired bootstrap resampling~\cite{efron1979bootstrap} is the standard way to attach confidence intervals to differences in evaluation metrics~\cite{koehn2004significance}, and Dror et al.~\cite{dror2018hitchhiker} give a protocol for choosing among significance tests in natural language processing. Card et al.~\cite{card2020power} showed that many comparisons in that field are underpowered, and Miller~\cite{miller2024} treats an evaluation as an experiment on questions drawn from a super-population, recommends reporting standard errors, clustered when questions are not independent, and gives formulas for comparing two models and for planning an evaluation's size. The Wilson score interval~\cite{wilson1927} keeps binomial intervals inside [0, 1] and performs well at small sample sizes~\cite{agresti1998}. The harness combines these elements: Wilson intervals for rates, a paired cluster bootstrap for the difference, an exact McNemar test on the discordant items, and a minimum detectable effect in every report. Its decision rule follows the logic of non-inferiority testing~\cite{walker2011}, with the margin, the tie policy and the optional requirement of a significant gain made explicit.

\subsection{Testing machine-learning systems}

Breck et al.~\cite{breck2017mltestscore} propose 28 specific tests and monitoring needs as a rubric for the production readiness of machine-learning systems, observing that such tests are hard to write because a model's behaviour is hard to specify in advance. CheckList~\cite{ribeiro2020checklist} brings behavioural testing from software engineering to natural language processing models. The harness applies the regression-testing part of this agenda to an LLM feature and runs it in CI. To make CI deterministic and free of model calls, it records model replies and replays them, in the manner of HTTP record-and-replay test libraries such as VCR.py~\cite{vcrpy}, but keys each recording to the SHA-256 of the full request, so that a changed prompt, model name or sampling setting invalidates the recording instead of silently reusing it.

\subsection{Adaptation, quantization and serving}

LoRA~\cite{hu2022lora} adapts a frozen model by training low-rank update matrices; with the update's second factor initialised to zero, an adapter is an exact identity, which Experiment A uses as a no-op control. Post-training quantization such as GPTQ~\cite{frantar2023gptq} reduces weights to 3 or 4 bits with little loss, and BitNet b1.58~\cite{ma2024bitnet} trains models whose weights are ternary, \{$-$1, 0, +1\}, from scratch. Ternary Bonsai 2~\cite{bonsai2} is a ternary model published with Qwen3.8-27B~\cite{qwen38card} as its base model and distributed as GGUF files for llama.cpp~\cite{llamacpp}; Experiment C compares it with the bf16 model. vLLM~\cite{kwon2023vllm} serves models with paged key-value memory and continuous batching, and speculative decoding~\cite{leviathan2023speculative} accelerates generation without changing the output distribution in exact arithmetic. In floating point, however, results depend on batch composition and reduction order: Yuan et al.~\cite{yuan2025nondeterminism} report accuracy differences of up to 9\% for a reasoning model under greedy decoding in bfloat16 as the GPU count, GPU type and batch size change, and He et al.~\cite{he2025nondeterminism} attribute most run-to-run variation of inference servers to kernels that are not invariant to batch size. Experiment A measures the consequence for a regression gate directly.

\section{System Design}

The harness is a Python package (\texttt{src/\allowbreak{}llm\_\allowbreak{}eval\_\allowbreak{}harness/\allowbreak{}}, standard library only, Python 3.10 or later) with a command-line interface, \texttt{llm-{}eval}. Its full specification is \texttt{docs/\allowbreak{}DESIGN.\allowbreak{}md}; the module docstrings state the same contracts in more detail. This section summarises the design decisions on which the experiments depend.

\subsection{Architecture and data flow}

An evaluation proceeds through the following stages, each owned by one module:

\begin{enumerate}
\item \textbf{Suite.} \texttt{tasks.\allowbreak{}load\_\allowbreak{}suite} reads a suite directory (\texttt{suite.\allowbreak{}json}, \texttt{tasks.\allowbreak{}jsonl}) and \texttt{tasks.\allowbreak{}split\_\allowbreak{}of} assigns every task to the development (dev) or held-out split. \texttt{validation.\allowbreak{}validate\_\allowbreak{}suite} checks the suite before anything is scored on it.
\item \textbf{Items.} A feature (\texttt{features.\allowbreak{}py}) turns the tasks of one split into items: \texttt{bugfix:<task>} for the bug-fix feature, and \texttt{judge:<task>:reference} and \texttt{judge:<task>:wrong} for the judge.
\item \textbf{Requests.} \texttt{prompts.\allowbreak{}build\_\allowbreak{}request} renders a versioned template for an item into a chat request. Its SHA-256, computed over the canonical JSON of exactly the body that the live client sends, identifies the request.
\item \textbf{Replies.} A live OpenAI-compatible client (\texttt{clients.\allowbreak{}OpenAICompatClient}, built on \texttt{urllib}) answers requests; a recording client stores each reply in a cassette; a replay client (\texttt{clients.\allowbreak{}ReplayClient}) answers from the cassette and refuses a recording whose request hash differs from the rebuilt request.
\item \textbf{Scoring.} The feature scores each reply (Section III-C) into an \texttt{ItemResult}.
\item \textbf{Metrics and comparison.} \texttt{metrics.\allowbreak{}arm\_\allowbreak{}metrics} aggregates one arm; \texttt{metrics.\allowbreak{}compare} joins two arms on item identifiers.
\item \textbf{Decision.} \texttt{gate.\allowbreak{}decide} applies the gate rule (Section III-G) and produces a JSON report, the only output written to standard output.
\item \textbf{Record keeping.} \texttt{ledger.\allowbreak{}py} appends evaluation, promotion and rollback rows to a hash-chained ledger.
\end{enumerate}

Only the shared records in \texttt{records.\allowbreak{}py} cross module boundaries, and \texttt{metrics} and \texttt{gate} depend on nothing but those records. \texttt{tests/\allowbreak{}test\_\allowbreak{}contract.\allowbreak{}py} fails if a module imports anything outside the standard library.

\subsection{The features under test}

\textbf{Bug fix.} The model receives a buggy Python function with its docstring and one test that fails on it, rendered as a single assertion. It must answer with an optional thinking block followed by exactly one JSON object with the keys \texttt{explanation} and \texttt{patched\_\allowbreak{}code}. A product would apply \texttt{patched\_\allowbreak{}code} directly, so the harness scores exactly that. The model never sees the task's hidden tests, its reference fix or its wrong fix. Templates are versioned (\texttt{bugfix/\allowbreak{}v1}, \texttt{bugfix/\allowbreak{}v2}) and never change once published; their SHA-256 values are pinned by \texttt{tests/\allowbreak{}test\_\allowbreak{}prompts.\allowbreak{}py}.

\textbf{Judge.} The model receives a task and a proposed patch and answers \texttt{\{"verdict": true\}} or \texttt{\{"verdict": false\}}. Each task yields two items: its reference fix (label true) and a plausible wrong fix that passes the shown test and fails at least one hidden test (label false). Labels therefore come from execution, never from another model. Because every task contributes one item of each class, a judge that decides from the task alone, without reading the patch, scores exactly 0.5 balanced accuracy.

\subsection{Scoring a reply}

A bug-fix reply is checked in a fixed order, and each check is recorded in the item's result:

\begin{enumerate}
\item \textbf{Thinking.} Thinking may arrive in a separate field of the server's response or as a \texttt{<think>} block at the start of the content. It is valid when it is closed (an answer follows it), lies within the feature's word limits (1 to 1,200 words for both features), and appears in only one place; it is optional for both features.
\item \textbf{Schema.} A reply that ended at the token limit (\texttt{finish\_\allowbreak{}reason} ``length'') is invalid. Otherwise the text after the thinking must be exactly one JSON object with exactly the required keys, of the required types, and non-blank strings: no surrounding prose, code fence, duplicate key or trailing text. A lenient parse that strips a fence or extracts the outermost braces is recorded as a diagnostic only and never makes a reply valid.
\item \textbf{Static scan.} The patch is scanned (Section III-D) and the findings it added relative to the buggy code are recorded.
\item \textbf{Execution.} Only a schema-valid patch without an added block-class finding is executed, against the shown test and then the hidden tests (Section III-E).
\end{enumerate}

An item is \emph{resolved} when its reply is schema-valid, adds no block-class finding, and passes every shown and hidden test. An item whose shown test passed and a hidden test failed is an \emph{overfit}. A judge reply's prediction is the verdict of a schema-valid reply; a reply without a valid verdict is wrong for either label, and no code runs.

\subsection{Static analysis}

\texttt{sast.\allowbreak{}scan} parses the patch and resolves imports, aliases, simple assignments and \texttt{getattr(module, "name")}, so that \texttt{from os import system as run; run(cmd)} is found; when the code does not parse, the block-class rules run as regular expressions on the raw text. A finding belongs to the patch when the buggy code did not already contain it (a multiset difference on rule and snippet). Six block-class rules make an item unresolved and prevent its code from running; six flag-class rules are counted and decide nothing.

\begin{table}[!t]
\caption{The block-class static-analysis rules (\texttt{src/\allowbreak{}llm\_\allowbreak{}eval\_\allowbreak{}harness/\allowbreak{}sast.\allowbreak{}py}). Flag-class rules (subprocess, dynamic import, file writes, network, introspection, \texttt{sys.\allowbreak{}exit}) are reported but decide nothing.}
\label{tab:sast}
\centering
\scriptsize
\setlength{\tabcolsep}{3pt}
\begin{tabular}{l l p{0.45\columnwidth}}
\hline
Rule & CWE & Matches \\
\hline
\texttt{eval-{}exec} & CWE-95 & calls of the built-in \texttt{eval} or \texttt{exec} \\
\texttt{os-{}command} & CWE-78 & \texttt{os.\allowbreak{}system}, \texttt{os.\allowbreak{}popen}, \texttt{os.\allowbreak{}exec*}, \texttt{os.\allowbreak{}spawn*}, \texttt{pty.\allowbreak{}spawn} \\
\texttt{shell-{}true} & CWE-78 & any call with \texttt{shell=True} \\
\texttt{unsafe-{}deserialization} & CWE-502 & \texttt{pickle}, \texttt{marshal} or \texttt{dill} loads, \texttt{shelve.\allowbreak{}open}, unsafe \texttt{yaml.\allowbreak{}load} \\
\texttt{process-{}control} &  & \texttt{os.\allowbreak{}\_\allowbreak{}exit}, \texttt{os.\allowbreak{}kill}, \texttt{atexit.\allowbreak{}register}, \texttt{signal.\allowbreak{}signal}, \texttt{signal.\allowbreak{}alarm} \\
\texttt{interpreter-{}tamper} &  & \texttt{sys.\allowbreak{}settrace}, \texttt{sys.\allowbreak{}setprofile}, writes into \texttt{sys.\allowbreak{}modules} or \texttt{builtins} \\
\hline
\end{tabular}
\end{table}

A block-class rule is admitted only if it never fires on a correct fix: suite validation checks every reference solution against it.

\subsection{Execution and the sandbox}

Tests are stored as data: a call expression and either the Python literal it must return, optionally with a numeric tolerance, or the exception class it must raise. The child process that runs model-written code receives only the call expressions and reports, per call, the \texttt{repr} of the returned value or the class names of the raised exception. The parent process decides the outcome (execution rule \texttt{values/\allowbreak{}1}). A test passes only when its record exists, is the only record for its index, and matches the expectation; a missing record is a failure, so an early \texttt{os.\allowbreak{}\_\allowbreak{}exit}, a crash or a skipped call cannot pass, and an exit status of 0 is never a pass by itself. If the sandbox itself cannot start, the evaluation stops with an infrastructure error instead of scoring the item as failed.

The sandbox runs the child as \texttt{python -{}I -{}S} in a fresh temporary directory that is its working directory, \texttt{HOME} and \texttt{TMPDIR}, with an environment built from nothing, under a deadline (10 s per item by default) after which the whole process group is killed, and with best-effort POSIX limits on CPU time, address space (1 GiB) and file size (10 MiB). It is explicitly not a security boundary: the child runs as the invoking user and can read that user's files and open network connections. In CI the only code executed is the code in committed cassettes, reviewed like any other file; untrusted recordings should be scored in a container without network access.

\subsection{Statistics}

\textbf{Rates.} Every rate $k/n$ carries a Wilson score interval~\cite{wilson1927} with $z = 1.96$:

\begin{equation}
\frac{k + z^2/2}{n + z^2} \pm \frac{z \sqrt{k(n-k)/n + z^2/4}}{n + z^2},
\end{equation}

which, unlike the normal approximation, stays inside [0, 1] and does not collapse at $k = 0$ or $k = n$~\cite{agresti1998}.

\textbf{Judge metrics.} With true positives and false negatives counted on reference items and true negatives and false positives on wrong items, and a reply without a valid verdict counted as wrong in its class, the primary metric is balanced accuracy,

\begin{equation}
\mathrm{BA} = \tfrac{1}{2}\left(\frac{TP}{P} + \frac{TN}{N}\right),
\end{equation}

with a bootstrap interval over tasks. A report flags a judge whose interval does not lie above 0.5.

\textbf{Paired comparison.} The arms are joined on item identifiers. For item $i$, let $x_i = c_i - b_i$ with $c_i, b_i \in \{0, 1\}$ the correctness of the candidate and of the baseline. \emph{Gains} are the items with $x_i = +1$, \emph{losses} those with $x_i = -1$, and $\Delta$ is the primary metric of the candidate minus that of the baseline on the joined items, $(\text{gains} - \text{losses})/n$ for the resolved rate. The 95\% interval of $\Delta$ is a percentile interval from 2,000 paired bootstrap draws~\cite{efron1979bootstrap} of \emph{clusters} with replacement: a task's two judge items, and twin tasks that share a group, are drawn together because they are not independent. The generator is seeded with the string \texttt{<seed>/\allowbreak{}<feature>/\allowbreak{}<split>}, so the interval is identical on every machine.

\textbf{McNemar's test.} With $b$ losses and $c$ gains, the exact two-sided McNemar p-value~\cite{mcnemar1947} is

\begin{equation}
p = \min\left(1,\; 2 \sum_{k=0}^{\min(b,c)} \binom{b+c}{k} 2^{-(b+c)}\right),
\end{equation}

and $p = 1$ when $b + c = 0$. With no gains, at least six losses are needed for $p < 0.05$ (since $2 \cdot 2^{-5} = 0.0625$ and $2 \cdot 2^{-6} = 0.03125$), whatever the suite size; every report states this number.

\textbf{Minimum detectable effect.} With $\widehat{\mathrm{se}}$ the standard deviation of the bootstrap deltas, the report gives $\mathrm{MDE}_{80} = (z_{0.975} + z_{0.80})\,\widehat{\mathrm{se}} \approx 2.80\,\widehat{\mathrm{se}}$: the smallest true change that a two-sided 5\% test on this suite would detect with 80\% power.

\subsection{The gate rule}

Rule \texttt{noninferior/\allowbreak{}1} promotes a candidate only when every applicable gate passes. Table~\ref{tab:gates} lists them.

\begin{table}[!t]
\caption{The gates of rule \texttt{noninferior/\allowbreak{}1} (\texttt{src/\allowbreak{}llm\_\allowbreak{}eval\_\allowbreak{}harness/\allowbreak{}gate.\allowbreak{}py}). The defaults are margin 0, ties pass, and alpha 0.05.}
\label{tab:gates}
\centering
\scriptsize
\setlength{\tabcolsep}{3pt}
\begin{tabular}{l p{0.60\columnwidth}}
\hline
Gate & Passes when \\
\hline
\texttt{schema\_\allowbreak{}validity} & the candidate's schema validity is at least 0.95 \\
\texttt{think\_\allowbreak{}validity} & the candidate's thinking validity is at least 0.90 \\
\texttt{new\_\allowbreak{}sast\_\allowbreak{}violations} & no item has a block-class finding in the candidate's reply whose rule is absent from the baseline's reply to the same item (bug fix only) \\
\texttt{primary} & $\Delta \geq -\text{margin}$ (ties pass), or $\Delta > -\text{margin}$ with \texttt{-{}-{}tie fail} \\
\texttt{regression\_\allowbreak{}test} & not (losses > gains and McNemar $p <$ alpha) \\
\texttt{gain\_\allowbreak{}test} & only with \texttt{-{}-{}require-{}gain}: gains > losses and McNemar $p <$ alpha \\
\hline
\end{tabular}
\end{table}

The tie policy is part of the rule because a tie is a decision. Passing ties is the right default for a regression test, which should not block a refactoring that changes nothing; replacing a baseline because a change is meant to help should demand more, \texttt{-{}-{}tie fail -{}-{}require-{}gain}. A positive margin absorbs the noise of a nondeterministic endpoint. The gate decides on the point estimate, with McNemar's test as a second guard, because on suites of a few dozen items the 95\% interval of $\Delta$ is wider than any margin worth setting: requiring the interval's lower end to exceed the negative margin would fail nearly every candidate, and requiring a significant regression before rejecting would pass nearly every one. Without a baseline, the primary metric is compared with a floor (0.5 for the bug fix, 0.7 for the judge).

The gate refuses, with exit status 2 and no report, to compare arms of different features, splits or suites, arms that do not cover the same item identifiers, and recordings that do not answer the requests rebuilt today. Exit status 0 means only that the evaluation ran; CI reads the verdict from the report (\texttt{jq -{}e '.\allowbreak{}promote == true'}), so a rejection and a failure to run remain distinguishable. The report contains no time stamps, host names or absolute paths.

\subsection{The ledger}

\texttt{ledger/\allowbreak{}ledger.\allowbreak{}jsonl} is committed with the code. Each row carries the SHA-256 of the previous line and a hash of its own content, so an edited, removed or reordered row is detected by \texttt{llm-{}eval ledger verify}, which also re-fingerprints each feature's current baseline directory. A promotion is refused unless its report is a passing held-out report of exactly the candidate being promoted against the current baseline; an override must state a reason and is recorded as one. A rollback pops the stack of promotions. Evaluation rows, written by \texttt{llm-{}eval gate -{}-{}record-{}evaluation}, count the distinct candidates judged on the held-out split; once ten have been judged, every report warns that the split is spent. CI reads the ledger and never writes it.

\subsection{Continuous integration}

\begin{table*}[!t]
\caption{Jobs of \texttt{.\allowbreak{}github/\allowbreak{}workflows/\allowbreak{}ci.\allowbreak{}yml}. The workflow has read-only repository permissions.}
\label{tab:ci}
\centering
\scriptsize
\setlength{\tabcolsep}{3pt}
\begin{tabular}{p{0.13\textwidth} p{0.22\textwidth} p{0.55\textwidth}}
\hline
Job & Trigger & Function \\
\hline
\texttt{tests} & every push to \texttt{main}, every pull request & unit and end-to-end tests on Python 3.10, 3.11 and 3.12; suite validation; ledger verification \\
\texttt{regression-{}gate} & after \texttt{tests} & replays the demo cassettes, gates the candidate and a deliberately regressed candidate against the current baseline on the held-out split, writes the job summary, uploads the reports, then requires the candidate to promote and the regressed candidate to be rejected by the expected gates \\
\texttt{experiments} & after \texttt{tests} & replays every experiment's recorded real-model cassettes and fails if the current harness does not reproduce its committed reports byte for byte \\
\texttt{live-{}config}, \texttt{live-{}eval} & manual dispatch & when repository secrets configure an endpoint, record a candidate and a baseline live and gate them \\
\hline
\end{tabular}
\end{table*}

The self-test in \texttt{regression-{}gate} makes a harness change that would let a clearly regressed candidate through fail the build even though the real candidate still passes. \texttt{tests/\allowbreak{}test\_\allowbreak{}workflow.\allowbreak{}py} parses the workflow without a YAML library, checks every \texttt{llm-{}eval} command line in it against the command-line parser, and runs the gate job's scripts, including their failure paths, on a copy of the repository.

\subsection{Report generation}

The result tables of this report are generated by \texttt{scripts/\allowbreak{}make\_\allowbreak{}report\_\allowbreak{}tables.\allowbreak{}py} from the committed gate reports in \texttt{experiments/\allowbreak{}*/\allowbreak{}expected/\allowbreak{}*.\allowbreak{}json}, as markdown (\texttt{docs/\allowbreak{}report/\allowbreak{}tables.\allowbreak{}md} and marked blocks in Section V) and as LaTeX (\texttt{docs/\allowbreak{}paper/\allowbreak{}tables/\allowbreak{}*.\allowbreak{}tex}). The script only formats; it computes no statistic. \texttt{tests/\allowbreak{}test\_\allowbreak{}report\_\allowbreak{}tables.\allowbreak{}py} fails when a committed table differs from the one the reports produce, so the tables cannot drift from the data that CI reproduces.

\section{Experimental Methodology}

\subsection{Task suites and splits}

\textbf{Demo suite.} \texttt{tasks/\allowbreak{}demo} holds 34 tasks written for the repository and taken from no dataset or code base (\texttt{"synthetic": true} in \texttt{suite.\allowbreak{}json}). Each task is a small pure Python function with one bug, tagged by bug class (\texttt{off-{}by-{}one}, \texttt{boundary} and others), one shown test that fails on the buggy code, at least two hidden tests, a reference fix and a plausible wrong fix. \texttt{llm-{}eval tasks validate}, run by CI on every Python version, checks in the sandbox that the buggy code fails its shown test, that the reference passes every test twice with the same outcomes, that the wrong fix passes the shown test and fails a hidden one, that no block-class rule fires on any of the three code fields, and that no rendered prompt contains a hidden test's call, the reference or the wrong fix.

\textbf{Split.} A task's split is a deterministic function of a stable key, never of its text:

\begin{equation}
\mathrm{bucket} = \mathrm{int}\big(\mathrm{sha256}(\mathrm{key})[0{:}8],\,16\big) \bmod 100,
\end{equation}

where the key is the task's normalised \texttt{group} or, without one, its identifier, and the task is held out when the bucket is below the suite's \texttt{heldout\_\allowbreak{}pct}. Editing a task's wording therefore never moves it across the split, and adding a task never moves another. Twins, tasks that pose the same problem reworded or with a renamed function, share a group and therefore a side; \texttt{tasks.\allowbreak{}find\_\allowbreak{}twins} warns about near-identical code across the split and validation fails on identical code. The demo suite sets \texttt{heldout\_\allowbreak{}pct} to 50 so that its held-out split is large enough to detect anything at all; the hash places 21 tasks (20 clusters, because \texttt{digit-{}sum} and \texttt{sum-{}of-{}digits} are twins) in the held-out split and 13 in the dev split. The bug-fix feature thus has 21 held-out and 13 dev items, and the judge 42 and 26.

\textbf{HumanEvalFix suite (Experiment C, planned).} Experiment C is designed to use the Python problems of HumanEvalFix~\cite{muennighoff2024octopack}, whose buggy functions derive from HumanEval~\cite{chen2021codex}, converted to the harness's suite format by \texttt{scripts/\allowbreak{}import\_\allowbreak{}humanevalfix.\allowbreak{}py} into \texttt{tasks/\allowbreak{}humanevalfix/\allowbreak{}}. The suite's \texttt{README.\allowbreak{}md} and \texttt{IMPORT.\allowbreak{}json} document each conversion choice and every task or assertion the format cannot express exactly and that is therefore left out. The suite's final size and split are to be reported with Experiment C's results, from the committed suite; no claim about them is made here.

\subsection{Arms}

\textbf{Controlled synthetic arms.} \texttt{cassettes/\allowbreak{}baseline}, \texttt{cassettes/\allowbreak{}candidate} and \texttt{cassettes/\allowbreak{}candidate-{}regressed} were generated by \texttt{scripts/\allowbreak{}make\_\allowbreak{}demo\_\allowbreak{}cassettes.\allowbreak{}py}, not by a model, from documented deterministic rules over each task's position in its split: the reference code stands for a fix, the wrong code for an overfitted patch, and hand-written \texttt{eval()} variants, prose around the JSON, a missing key and replies truncated at the token limit stand for the failures. The baseline uses prompts \texttt{bugfix/\allowbreak{}v1} and \texttt{judge/\allowbreak{}v1}; the candidate uses \texttt{v2}, which adds the instruction to fix the function for every input its docstring describes and not to special-case test inputs; the regressed candidate uses \texttt{v2} with a nominally smaller model. These arms have known ground truth by construction and test the gate, not any model.

\textbf{Experiment A: LoRA adapters on an open-weight model} (\texttt{experiments/\allowbreak{}exp-{}a-{}qwen38/\allowbreak{}}). The base model is Qwen3.8-27B~\cite{qwen38card}, in a community fine-tune whose refusal behaviour was removed by ``abliteration'', a weight edit along the refusal direction identified by Arditi et al.~\cite{arditi2024refusal}, in bf16, served by vLLM~\cite{kwon2023vllm} on two A100 GPUs with thinking enabled (the series' thinking mode is described in~\cite{qwen3report}). Two arms were recorded on the same server:

\begin{itemize}
\item \texttt{qwen38-{}base}: the base model for both features (served model name \texttt{qwen38}).
\item \texttt{qwen38-{}adapters}: for the bug fix, an \emph{identity} LoRA adapter~\cite{hu2022lora} of rank 64 whose \texttt{lora\_\allowbreak{}B} matrices are all zero, so that the adapted model equals the base model in exact arithmetic (served name \texttt{coder-{}lora}); for the judge, a \emph{trained} LoRA adapter of rank 64, trained on the base model's own rejection-sampled verdicts on public code-review data (served name \texttt{judge-{}lora}).
\end{itemize}

Both arms used prompts \texttt{bugfix/\allowbreak{}v2} and \texttt{judge/\allowbreak{}v2}. The identity adapter is a no-op control: any difference it produces measures the serving stack, not the adapter.

\textbf{Experiment C: weight quantization and serving engine (planned).} Three arms are designed for the bug-fix feature on the HumanEvalFix suite: Qwen3.8-27B in bf16 served by vLLM; Ternary Bonsai 2 27B~\cite{bonsai2}, a ternary model with Qwen3.8-27B as its base, as a Q2\_0 GGUF file served by llama.cpp~\cite{llamacpp}; and Qwen3.8-27B as a Q8\_0 GGUF file served by llama.cpp, the engine control. Three paired comparisons follow: C1, Q8\_0 against bf16, measures the change of engine (together with 8-bit rounding, expected to be small); C2, ternary against Q8\_0, measures the ternary model at a fixed engine; and C3, ternary against bf16, measures the whole deployment change, as a product team would experience it. Without the control arm, C3 alone could not separate the effect of the weights from that of the engine.

\subsection{Recording protocol}

Replies are recorded with \texttt{llm-{}eval record}, which builds each item's request exactly as replay will rebuild it and stores the model name, prompt identifier and settings in the cassette directory's \texttt{manifest.\allowbreak{}json}, and each reply under its request's SHA-256. Experiment A used temperature 0, seed 0, at most 16,384 completion tokens and \texttt{\{"chat\_\allowbreak{}template\_\allowbreak{}kwargs": \{"enable\_\allowbreak{}thinking": true\}\}} (\texttt{experiments/\allowbreak{}exp-{}a-{}qwen38/\allowbreak{}*/\allowbreak{}manifest.\allowbreak{}json}). The server returned the thinking in a separate reasoning field, which the harness scores like a \texttt{<think>} block. A request that fails after the client's retries is not recorded, so an unreliable endpoint cannot be scored as a bad model; re-running the command resumes. Recording into a directory whose recordings of the same feature used a different model, prompt, settings or suite version is refused. Each arm holds one reply per item: the experiments sample once at temperature 0.

\subsection{Replay determinism}

Replay rebuilds every request from the manifest and the suite and compares its SHA-256 with the recorded one; a mismatch or a missing recording refuses the evaluation. The bootstrap is seeded per feature and split, and the report is serialised canonically without time stamps, host names or absolute paths. Replaying the same cassettes therefore yields the same bytes on any machine. The committed reports of Experiment A are in \texttt{experiments/\allowbreak{}exp-{}a-{}qwen38/\allowbreak{}expected/\allowbreak{}}, one per feature and split, produced by

{\footnotesize
\begin{verbatim}
llm-eval gate --feature <f> --split <s> \
  --candidate \
    experiments/exp-a-qwen38/qwen38-adapters \
  --baseline \
    experiments/exp-a-qwen38/qwen38-base
\end{verbatim}
}

with the default gate configuration (margin 0, ties pass, schema validity at least 0.95, thinking validity at least 0.90, alpha 0.05, 2,000 bootstrap draws, seed 0). The CI job \texttt{experiments} reruns these commands and compares the output with the committed files using \texttt{cmp}.

\subsection{Hypotheses}

The following hypotheses state what the design of each arm predicts. They are listed before any result. H1 and H2 follow from the harness's specification; H3 and H4 state what the adapters were expected to do, and were written down in this form after Experiment A's replies had been recorded (Section VI-E). H5 and H6 are fixed before Experiment C is recorded.

\begin{itemize}
\item \textbf{H1 (gate validity on controlled arms).} On the held-out demo split, the synthetic candidate promotes for both features, and the synthetic regressed candidate is rejected, failing at least \texttt{primary}, \texttt{schema\_\allowbreak{}validity} and \texttt{new\_\allowbreak{}sast\_\allowbreak{}violations} for the bug fix and \texttt{primary} and \texttt{schema\_\allowbreak{}validity} for the judge.
\item \textbf{H2 (replay determinism).} Replaying the committed cassettes with the committed harness reproduces every committed report byte for byte.
\item \textbf{H3 (no-op invariance).} The identity adapter's bug-fix replies equal the base model's on every item, so that $\Delta = 0$ with no gains and no losses on both splits, and the gate promotes under the default tie policy.
\item \textbf{H4 (trained judge).} The trained judge adapter raises balanced accuracy over the base model ($\Delta > 0$) on both splits and is promoted on both.
\item \textbf{H5 (engine control, Experiment C).} Comparison C1 (Q8\_0 under llama.cpp against bf16 under vLLM) is not rejected by \texttt{regression\_\allowbreak{}test}, and its number of discordant items is of the same order as the no-op discordance measured in Experiment A.
\item \textbf{H6 (ternary weights, Experiment C).} Comparison C2 (ternary against Q8\_0, both under llama.cpp) is promoted by rule \texttt{noninferior/\allowbreak{}1} with a margin fixed before recording.
\end{itemize}

\subsection{Analysis plan}

Each report is one decision and is interpreted on its own; no correction for multiple comparisons is applied, and Section VI discusses the consequence. Every decision is read together with its interval, $\mathrm{MDE}_{80}$ and the number of discordant items. For every discordant item, the per-item rows written by \texttt{llm-{}eval gate -{}-{}items-{}out DIR} are inspected to classify the disagreement as a functional failure (a test failed or a verdict was wrong) or a format failure (truncation, invalid thinking, or a reply that is not a strictly valid JSON object). Cost is reported as mean completion tokens per item and median latency, and is not gated.

\section{Results}

The tables in this section, except the demo table in Section V-A, are generated by \texttt{scripts/\allowbreak{}make\_\allowbreak{}report\_\allowbreak{}tables.\allowbreak{}py} from the committed gate reports and must not be edited by hand.

\subsection{Controlled arms (H1)}

\begin{table*}[!t]
\caption{Gate decisions on the held-out demo split for the synthetic arms (21 tasks; 21 bug-fix and 42 judge items), against the ledger's current baseline \texttt{cassettes/\allowbreak{}baseline}. Every arm is synthetic. Source: \texttt{README.\allowbreak{}md} (``The demo results''), pinned by \texttt{tests/\allowbreak{}test\_\allowbreak{}demo.\allowbreak{}py}.}
\label{tab:demo}
\centering
\scriptsize
\setlength{\tabcolsep}{3pt}
\begin{tabular}{lllllllp{0.22\textwidth}}
\hline
Feature & Arm & Primary & $\Delta$ [95\% CI] & Gains / losses & McNemar $p$ & Verdict & Failed gates \\
\hline
bug fix & baseline (\texttt{bugfix/\allowbreak{}v1}) & 14/21 = 0.667 &  &  &  & current &  \\
bug fix & candidate (\texttt{bugfix/\allowbreak{}v2}) & 17/21 = 0.810 & +0.143 [$-$0.045, +0.350] & 4 / 1 & 0.375 & PROMOTE & none \\
bug fix & candidate-regressed & 10/21 = 0.476 & $-$0.190 [$-$0.391, +0.000] & 0 / 4 & 0.125 & REJECT & primary, schema\_validity, new\_sast\_violations \\
judge & baseline (\texttt{judge/\allowbreak{}v1}) & 0.571 &  &  &  & current &  \\
judge & candidate (\texttt{judge/\allowbreak{}v2}) & 0.881 & +0.310 [+0.175, +0.429] & 15 / 2 & 0.0023 & PROMOTE & none \\
judge & candidate-regressed & 0.429 & $-$0.143 [$-$0.350, +0.048] & 6 / 12 & 0.238 & REJECT & primary, schema\_validity \\
\hline
\end{tabular}
\end{table*}

H1 holds: both candidates promote, and both regressed candidates are rejected by exactly the predicted gates. The same verdicts hold on the dev split. Two properties of the decisions matter for what follows. The bug-fix candidate promotes because it is not worse, not because its improvement is significant: its interval includes 0 and McNemar's p is 0.375, and with \texttt{-{}-{}tie fail -{}-{}require-{}gain} it is rejected on \texttt{gain\_\allowbreak{}test} while the judge candidate (p = 0.0023) still promotes. The regressed bug-fix candidate's four losses are not significant (p = 0.125, below the six losses needed), so it is rejected by the point estimate, by its invalid replies (schema validity 0.762) and by the \texttt{eval()} call it introduced, a new block-class finding whose patch was never executed. The baseline judge approves 39 of its 42 patches; its balanced accuracy of 0.571 has the interval [0.500, 0.650], and the report flags it as indistinguishable from chance.

\subsection{Replay determinism (H2)}

H2 holds for the committed tree. The CI job \texttt{experiments} reruns the four gate commands of Section IV-D and compares each output with the committed report using \texttt{cmp}; replaying the four reports with the harness at the commit of this report produced byte-identical files. For the synthetic arms, \texttt{scripts/\allowbreak{}make\_\allowbreak{}demo\_\allowbreak{}cassettes.\allowbreak{}py -{}-{}check}, run by \texttt{tests/\allowbreak{}test\_\allowbreak{}demo.\allowbreak{}py}, regenerates the three cassette sets and compares them byte for byte with the committed ones, and \texttt{llm-{}eval ledger verify} in the \texttt{tests} job re-fingerprints the current baseline.

\subsection{Experiment A: LoRA adapters on Qwen3.8-27B (H3, H4)}

Table~\ref{tab:exp-a-gate} gives the four gate decisions and Table~\ref{tab:exp-a-arms} the per-arm metrics.

% Generated by scripts/make_report_tables.py; do not edit. Regenerate with: python scripts/make_report_tables.py
\begin{table*}[!t]
\caption{Experiment A: gate decisions, candidate \texttt{qwen38-adapters} (model coder-lora, judge-lora) against baseline \texttt{qwen38-base} (model qwen38), rule \texttt{noninferior/1}, suite \texttt{demo}. Primary metric: resolved rate (bugfix), usable patch rate (patch), or balanced accuracy (judge, review). $\Delta$ is candidate minus baseline on paired items with a paired cluster-bootstrap 95\% interval; G/L are discordant gains and losses; MDE$_{80}$ is the minimum detectable effect at 80\% power. Generated from the committed reports in \texttt{experiments/exp-a-qwen38/expected}.}
\label{tab:exp-a-gate}
\centering
\scriptsize
\setlength{\tabcolsep}{2pt}
\begin{tabular}{llrllllrrlp{0.22\textwidth}}
\hline
Feature & Split & $n$ & Base & Cand. & $\Delta$ [95\% CI] & G/L & McNemar $p$ & MDE$_{80}$ & Verdict & Failed gates \\
\hline
bugfix & heldout & 21 & 1.000 (21/21) & 0.952 (20/21) & $-$0.048 [$-$0.150, +0.000] & 0 / 1 & 1.000 & 0.131 & REJECT & primary \\
bugfix & dev & 13 & 0.923 (12/13) & 0.846 (11/13) & $-$0.077 [$-$0.231, +0.000] & 0 / 1 & 1.000 & 0.210 & REJECT & schema\_validity (0.846), think\_validity (0.846), primary \\
judge & heldout & 42 & 1.000 & 0.976 & $-$0.024 [$-$0.075, +0.000] & 0 / 1 & 1.000 & 0.064 & REJECT & primary \\
judge & dev & 26 & 0.808 & 0.962 & +0.154 [+0.038, +0.269] & 4 / 0 & 0.125 & 0.177 & PROMOTE & none \\
\hline
\end{tabular}
\end{table*}


% Generated by scripts/make_report_tables.py; do not edit. Regenerate with: python scripts/make_report_tables.py
\begin{table*}[!t]
\caption{Experiment A: per-arm metrics. Intervals are Wilson (resolved rate) or a task-cluster bootstrap (balanced accuracy). Trunc.\ counts replies cut at the token limit; Tokens is the mean completion length.}
\label{tab:exp-a-arms}
\centering
\scriptsize
\setlength{\tabcolsep}{3pt}
\begin{tabular}{lllllrrrrrl}
\hline
Feature & Split & Arm & Model & Primary [95\% CI] & Schema & Think & Len.\ only & Trunc. & Tokens & Feature-specific \\
\hline
bugfix & heldout & qwen38-base & qwen38 & 1.000 [0.845, 1.000] & 1.000 & 1.000 & 0.000 & 0/21 & 412 & overfit 0.000, hidden pass 1.000, block SAST 0 \\
bugfix & heldout & qwen38-adapters & coder-lora & 0.952 [0.773, 0.992] & 0.952 & 0.952 & 0.000 & 0/21 & 477 & overfit 0.000, hidden pass 0.952, block SAST 0 \\
bugfix & dev & qwen38-base & qwen38 & 0.923 [0.667, 0.986] & 0.923 & 0.923 & 0.000 & 1/13 & 1852 & overfit 0.000, hidden pass 0.923, block SAST 0 \\
bugfix & dev & qwen38-adapters & coder-lora & 0.846 [0.578, 0.957] & 0.846 & 0.846 & 0.077 & 1/13 & 2625 & overfit 0.000, hidden pass 0.846, block SAST 0 \\
judge & heldout & qwen38-base & qwen38 & 1.000 [1.000, 1.000] & 1.000 & 1.000 & 0.000 & 0/42 & 395 & TPR 1.000, TNR 1.000, approve 0.500 \\
judge & heldout & qwen38-adapters & judge-lora & 0.976 [0.925, 1.000] & 0.976 & 0.952 & 0.000 & 1/42 & 899 & TPR 1.000, TNR 0.952, approve 0.500 \\
judge & dev & qwen38-base & qwen38 & 0.808 [0.691, 0.923] & 0.846 & 0.808 & 0.000 & 0/26 & 1784 & TPR 0.615, TNR 1.000, approve 0.308 \\
judge & dev & qwen38-adapters & judge-lora & 0.962 [0.885, 1.000] & 0.962 & 0.962 & 0.000 & 0/26 & 766 & TPR 0.923, TNR 1.000, approve 0.462 \\
\hline
\end{tabular}
\end{table*}


\textbf{H3 is refuted.} The identity adapter, which cannot change an output in exact arithmetic, produced one losing reply on each bug-fix split: resolved rate 0.952 against 1.000 on the held-out split and 0.846 against 0.923 on the dev split, with 0 gains and 1 loss each (McNemar p = 1.000). With margin 0 the \texttt{primary} gate fails and the adapter is rejected on both splits; on the dev split the single changed reply also lowers schema and thinking validity from the base model's 0.923 to 0.846. Both values are below the bounds of 0.95 and 0.90, so on this split the validity gates would reject the base model itself. The comparison itself marks the loss as noise: p = 1.000, and the report's warning states that at least six losses with no gains would be needed for significance.

\textbf{H4 is supported on the dev split only.} On the dev split the trained judge adapter raised balanced accuracy from 0.808 to 0.962 ($\Delta$ = +0.154, interval [+0.038, +0.269]) with 4 gains and 0 losses, and passed every gate. The gain is not significant by McNemar's test (p = 0.125): four discordant items in one direction cannot reach p < 0.05. On the held-out split the base model already scored 1.000, so the adapter could only tie or lose; one changed reply made it lose ($\Delta$ = $-$0.024, 0 gains, 1 loss, p = 1.000), and it was rejected on \texttt{primary}.

\textbf{Sensitivity.} The reports' minimum detectable effects at 80\% power are 0.131 (bug fix, held-out), 0.210 (bug fix, dev), 0.064 (judge, held-out) and 0.177 (judge, dev). On the held-out bug-fix split, a true change smaller than about three of 21 items is therefore unlikely to be detected.

\subsection{Item-level analysis}

The per-item rows (\texttt{-{}-{}items-{}out}) of the four reports show seven discordant items in total and classify them as follows:

\begin{itemize}
\item \textbf{Bug fix, held-out} (1 loss): the candidate's reply to \texttt{longest-{}run} placed thinking both in the reasoning field and in the content, and its answer was not a strictly valid JSON object, so the patch was never executed.
\item \textbf{Bug fix, dev} (1 loss): the candidate's reply to \texttt{parse-{}ratio} thought for 3,015 words, outside the 1 to 1,200 word window, and its answer was valid only under the lenient parse. A second dev item, \texttt{parse-{}int-{}list}, failed in both arms, each truncated at 16,384 completion tokens without an answer; it is concordant and does not enter the comparison.
\item \textbf{Judge, held-out} (1 loss): the candidate's reply to \texttt{average-{}rating:wrong} was truncated at 16,384 completion tokens without a verdict.
\item \textbf{Judge, dev} (4 gains): three of the base model's failures were format failures (on \texttt{first-{}repeat}, 1,225 words of thinking followed by no valid JSON object; on \texttt{parse-{}int-{}list} and \texttt{parse-{}ratio}, thinking that never closed), and one was a wrong verdict (\texttt{word-{}counts}, the reference patch judged incorrect after 7,227 words of thinking). All four involve reference items. \texttt{unique-{}in-{}order:reference} failed in both arms.
\end{itemize}

Six of the seven discordant items are therefore format or truncation failures, and one is a functional error. No bug-fix item in either arm on either split failed a test after a valid reply: the overfit rate is 0.000 throughout, and each arm's hidden-test pass rate equals its resolved rate.

\subsection{Cost}

Mean completion tokens per item are reported and not gated. The judge adapter used fewer tokens than the base model on the dev split (766 against 1,784), and more on the held-out split (899 against 395), where the report warns of a 2.3-fold increase; a single reply truncated at 16,384 tokens contributes 390 tokens to that held-out mean. The identity adapter's bug-fix replies averaged 477 against 412 tokens on the held-out split and 2,625 against 1,852 on the dev split, where each arm's mean includes one reply truncated at 16,384 tokens (1,260 tokens of a 13-item mean).

\subsection{Experiment C: quantization and serving engine (H5, H6)}

Experiment C has not been recorded. The tables below are placeholders that list the planned arms and comparisons; \texttt{scripts/\allowbreak{}make\_\allowbreak{}report\_\allowbreak{}tables.\allowbreak{}py} replaces them with the measured tables, under the same names, once \texttt{experiments/\allowbreak{}exp-{}c-{}*/\allowbreak{}expected/\allowbreak{}*.\allowbreak{}json} is committed. H5 and H6 remain untested.

% Generated by scripts/make_report_tables.py; do not edit. Regenerate with: python scripts/make_report_tables.py
\begin{table*}[!t]
\caption{Experiment C (planned): Ternary Bonsai 2 against bf16 Qwen3.8-27B on HumanEvalFix, gate decisions. \textbf{PLACEHOLDER: Experiment C has not been recorded; all numbers are pending. The table is regenerated from \texttt{experiments/exp-c-*/expected} when it exists.}}
\label{tab:exp-c-gate}
\centering
\scriptsize
\setlength{\tabcolsep}{3pt}
\begin{tabular}{lllp{0.25\textwidth}llll}
\hline
Comp. & Baseline arm & Candidate arm & Isolates & Base & Cand. & Delta [95\% CI] & Verdict \\
\hline
C1 & bf16-vllm & q8-llamacpp & the engine change (and 8-bit rounding): the engine control & pending & pending & pending & pending \\
C2 & q8-llamacpp & ternary-q2-llamacpp & the ternary model at a fixed engine & pending & pending & pending & pending \\
C3 & bf16-vllm & ternary-q2-llamacpp & the whole deployment change & pending & pending & pending & pending \\
\hline
\end{tabular}
\end{table*}


% Generated by scripts/make_report_tables.py; do not edit. Regenerate with: python scripts/make_report_tables.py
\begin{table*}[!t]
\caption{Experiment C: per-arm metrics. Arms: bf16-vllm: Qwen3.8-27B, bf16 weights, served by vLLM; bf16-vllm-rerun: bf16-vllm recorded a second time (the no-op arm); q8-llamacpp: Qwen3.8-27B, Q8\_0 GGUF, served by llama.cpp b12682 (engine control); ternary-llamacpp: Ternary Bonsai 2 27B, Q2\_0 GGUF, served by llama.cpp b12682. Intervals are Wilson (resolved rate, usable patch rate) or a task-cluster bootstrap (balanced accuracy). Trunc.\ counts replies cut at the token limit; Tokens is the mean and Total the sum of completion tokens.}
\label{tab:exp-c-arms}
\centering
\scriptsize
\setlength{\tabcolsep}{2pt}
\begin{tabular}{llllrrrrrrl}
\hline
Feature & Split & Arm & Primary [95\% CI] & Schema & Think & Len.\ only & Trunc. & Tokens & Total & Feature-specific \\
\hline
bugfix & heldout & bf16-vllm & 0.986 [0.927, 0.998] & 1.000 & 0.973 & 0.000 & 0/74 & 768 & 56818 & overfit 0.014, hidden pass 0.986, block SAST 0 \\
bugfix & heldout & bf16-vllm-rerun & 0.986 [0.927, 0.998] & 0.986 & 0.986 & 0.000 & 0/74 & 911 & 67379 & overfit 0.000, hidden pass 0.986, block SAST 0 \\
bugfix & heldout & q8-llamacpp & 1.000 [0.951, 1.000] & 1.000 & 1.000 & 0.000 & 0/74 & 722 & 53416 & overfit 0.000, hidden pass 1.000, block SAST 0 \\
bugfix & heldout & ternary-llamacpp & 0.959 [0.887, 0.986] & 0.986 & 0.959 & 0.000 & 1/74 & 1232 & 91141 & overfit 0.000, hidden pass 0.959, block SAST 0 \\
bugfix & dev & bf16-vllm & 1.000 [0.914, 1.000] & 1.000 & 1.000 & 0.000 & 0/41 & 703 & 28830 & overfit 0.000, hidden pass 1.000, block SAST 0 \\
bugfix & dev & bf16-vllm-rerun & 0.976 [0.874, 0.996] & 1.000 & 0.976 & 0.000 & 0/41 & 713 & 29215 & overfit 0.000, hidden pass 0.976, block SAST 0 \\
bugfix & dev & q8-llamacpp & 1.000 [0.914, 1.000] & 1.000 & 1.000 & 0.000 & 0/41 & 706 & 28961 & overfit 0.000, hidden pass 1.000, block SAST 0 \\
bugfix & dev & ternary-llamacpp & 0.976 [0.874, 0.996] & 0.976 & 1.000 & 0.000 & 0/41 & 881 & 36113 & overfit 0.000, hidden pass 0.976, block SAST 0 \\
judge & heldout & bf16-vllm & 0.953 [0.917, 0.981] & 0.973 & 0.973 & 0.020 & 0/148 & 744 & 110104 & TPR 0.905, TNR 1.000, approve 0.453 \\
judge & heldout & bf16-vllm-rerun & 0.973 [0.944, 0.993] & 0.993 & 0.959 & 0.007 & 0/148 & 743 & 109925 & TPR 0.946, TNR 1.000, approve 0.473 \\
judge & heldout & q8-llamacpp & 0.966 [0.938, 0.993] & 0.993 & 0.980 & 0.007 & 0/148 & 714 & 105708 & TPR 0.946, TNR 0.986, approve 0.473 \\
judge & heldout & ternary-llamacpp & 0.973 [0.944, 0.993] & 0.993 & 0.959 & 0.000 & 1/148 & 1004 & 148621 & TPR 0.946, TNR 1.000, approve 0.473 \\
judge & dev & bf16-vllm & 0.988 [0.963, 1.000] & 1.000 & 0.988 & 0.000 & 0/82 & 676 & 55420 & TPR 0.976, TNR 1.000, approve 0.488 \\
judge & dev & bf16-vllm-rerun & 0.963 [0.919, 1.000] & 0.988 & 0.988 & 0.012 & 0/82 & 670 & 54918 & TPR 0.927, TNR 1.000, approve 0.463 \\
judge & dev & q8-llamacpp & 1.000 [1.000, 1.000] & 1.000 & 0.976 & 0.000 & 0/82 & 673 & 55163 & TPR 1.000, TNR 1.000, approve 0.500 \\
judge & dev & ternary-llamacpp & 1.000 [1.000, 1.000] & 1.000 & 0.976 & 0.000 & 0/82 & 800 & 65613 & TPR 1.000, TNR 1.000, approve 0.500 \\
\hline
\end{tabular}
\end{table*}


\section{Discussion and Threats to Validity}

\subsection{What the decisions mean}

The controlled arms show that the gate does what its rule says: it promotes a candidate that is not worse, rejects one that is, and names the gates responsible (H1), and CI reproduces every decision exactly (H2). Experiment A shows what the rule means on a real model at small suite sizes. A margin of 0 rejected an adapter that changes nothing, on the evidence of one reply per split, and the significance guard correctly did not fire. Conversely, a real dev-split gain of four items was promoted but could not be confirmed, and the held-out split could not have confirmed any gain because the base model was at ceiling. Both outcomes are visible in the reports themselves (McNemar's p, $\mathrm{MDE}_{80}$ and the warning on the number of losses needed for significance), which is the purpose of reporting them with every decision.

Two practical rules follow. First, before a real comparison, a no-op arm should be recorded against the baseline to measure the serving noise floor, and the margin should be set from it, or promotion should additionally require that a loss be significant; Experiment A puts that floor at one discordant item in 13 to 21 on this suite and stack. Second, a held-out split on which the baseline scores at or near 1.000 cannot rank candidates, and a gain observed elsewhere should not be reported as confirmed until a harder or larger held-out suite is available.

\subsection{Internal validity}

\textbf{Serving nondeterminism.} The identity adapter is exactly equal to the base model in exact arithmetic, yet it changed replies. The serving stack in Experiment A applied the adapter through LoRA kernels, batched requests and used speculative decoding (\texttt{experiments/\allowbreak{}exp-{}a-{}qwen38/\allowbreak{}README.\allowbreak{}md}); each can change floating-point reduction order, and greedy decoding amplifies a single changed token into a different reply~\cite{yuan2025nondeterminism,he2025nondeterminism}. The experiment cannot attribute the change to one of these mechanisms, because it recorded one sample per item and did not vary them. Every difference between arms in this report therefore includes a nondeterministic component of unknown size beyond the one measured by the no-op arm.

\textbf{One sample per item.} Each arm holds one reply per item at temperature 0. A second recording of the same arm could differ on some items, and the paired analysis treats the recorded reply as the arm's fixed answer. The bootstrap interval reflects sampling of tasks, not sampling of replies.

\textbf{Engine confounds (Experiment C).} vLLM and llama.cpp differ in kernels, batching, sampling implementation, chat-template rendering and the parsing of reasoning output. A comparison of a ternary model under llama.cpp with a bf16 model under vLLM confounds the weights with the engine; the Q8\_0 arm under llama.cpp is designed to separate the two (C1 and C2). The separation is partial: C1 also includes 8-bit rounding, and C2 compares two different weight sets, because Ternary Bonsai 2 is published as its own model built from Qwen3.8-27B~\cite{bonsai2} rather than as a rounding of the same weights, so C2 measures the published ternary model and not ternary rounding as such.

\textbf{Recording conditions.} Both arms of Experiment A were recorded on the same server and with identical settings, so differences in hardware and settings do not enter the comparison; the order and load of requests during recording were not controlled.

\subsection{External validity}

The demo suite consists of 34 small, synthetic, single-function Python tasks, and Experiment A evaluates one model family, one base checkpoint (a community fine-tune rather than the stock release) and two adapters. Real bug fixing spans files and repositories with their own environments~\cite{jimenez2024swebench}; the harness's structure (paired items, execution labels, the gate and the ledger) carries over, but its sandbox does not, and per-task containers would be needed. HumanEvalFix, the planned suite of Experiment C, is also single-function Python. The findings about the gate's behaviour on small suites are expected to generalise; the absolute scores are not.

\subsection{Construct validity}

\textbf{Format and capability are confounded in the primary metric.} Six of the seven discordant items in Experiment A are format or truncation failures (Section V-D), and no bug-fix reply that was valid failed a test. The resolved rate and the judge's balanced accuracy, as defined, measure whether the feature produces a usable answer, which is what a product receives, but a reader interpreting them as reasoning ability would be misled. A per-item classification of failures, as in Section V-D, should accompany any interpretation.

\textbf{The thinking window.} The 1 to 1,200 word window and the 16,384-token limit turn long reasoning into failures. They encode a product decision about latency and cost, not a property of correctness; another product could set them differently and obtain different verdicts.

\textbf{Test adequacy.} ``Resolved'' means that the shown and hidden tests passed. Hidden tests reduce, but do not eliminate, the acceptance of wrong patches~\cite{liu2023evalplus}.

\textbf{The judge's task.} The judge is evaluated on exactly two patches per task, the reference and one wrong fix that passes the shown test. This design isolates whether the judge reads the patch, but it does not cover other kinds of wrong patch, and its labels are only as good as the tests.

\textbf{Static analysis.} The scan resolves common aliasing but cannot see a name constructed at run time, so the absence of a block-class finding is evidence, not proof, of the absence of the construct.

\subsection{Conclusion validity}

\textbf{Small suites and low power.} With 13 to 42 items per comparison, $\mathrm{MDE}_{80}$ lies between 0.064 and 0.210, and McNemar's test needs at least six losses with no gains to call a regression significant~\cite{card2020power}. A promotion is therefore evidence of non-inferiority at the stated margin, not of equality, and a rejection by one item is not evidence of a regression.

\textbf{Interval and test can disagree.} On the judge dev split the bootstrap interval of $\Delta$, [+0.038, +0.269], excludes 0 while McNemar's exact p is 0.125. With four discordant items, all in one direction, the exact test cannot reach significance, whereas a percentile bootstrap over only 13 clusters, whose resamples contain no loss at all, can understate the uncertainty at such sizes. The gate decides on neither alone; the report gives both, and neither should be cited without the other.

\textbf{Ceiling effects.} When the baseline scores 1.000, $\Delta \le 0$ is certain and any change can only tie or lose. Two of the four comparisons in Experiment A have a baseline at ceiling.

\textbf{Multiple comparisons.} Experiment A comprises four decisions and Experiment C will add at least three, without correction. Each decision is reported as a separate gate verdict and not as a family-wide claim.

\textbf{Post hoc hypotheses and reuse of the held-out split.} H3 and H4 were written in their present form after Experiment A's replies existed. They are the predictions implied by the arms' construction (an exact identity and a trained adapter), but they were not registered in advance. The demo suite's held-out split has been used for the synthetic arms and for Experiment A; the ledger counts held-out evaluations and warns after ten distinct candidates, a heuristic budget rather than a guarantee of the kind studied by Dwork et al.~\cite{dwork2015reusable}.

\textbf{Contamination.} The demo suite is public in the repository, so a model trained on the repository could have seen its held-out split; its held-out status protects against prompt tuning, not against training on the internet. HumanEvalFix derives from HumanEval, published in 2021 and widely copied, and may be in the pretraining data of the models compared in Experiment C~\cite{sainz2023contamination,jain2025livecodebench}. Contamination raises absolute scores and could interact with quantization if memorised solutions survive weight reduction differently from reasoned ones; a paired comparison removes neither effect.

\section{Conclusion and Future Work}

This report presented an evaluation harness that turns a change to an LLM feature into a paired, reproducible and explicitly specified promotion decision, enforced in CI without model calls. Scoring is grounded in execution, with hidden tests, values compared outside the process that runs the model's code, and a static scan that keeps unsafe patches from running. The gate applies a non-inferiority rule with an explicit tie policy and reports, with every decision, how large a change its suite could detect. Recorded replies make every report reproducible byte for byte, and a deliberately regressed candidate tests the gate in every CI run.

The study on a 27-billion-parameter model shows that these properties are necessary rather than ornamental. An identity adapter was rejected by a margin-0 gate on the strength of one reply per split, a dev-split gain of +0.154 balanced accuracy could not be confirmed on a held-out split at ceiling, and six of the seven disagreements between arms were failures of format or truncation rather than of reasoning. At the suite sizes a product team can afford, the decision rule, its margin and the suite's difficulty determine the verdict as much as the change under test does.

Future work follows directly from the threats in Section VI:

\begin{enumerate}
\item \textbf{Calibrate the margin from a no-op arm.} Record a no-op comparison for every serving stack and derive the margin, or a significance requirement on losses, from the measured discordance.
\item \textbf{Separate format from capability.} Report functional failures and format failures as separate metrics, so that a regression in output format is not read as a regression in reasoning, and vice versa.
\item \textbf{Harder, larger and private suites.} Replace the ceiling-bound demo split with a suite on which the baseline is far from 1.000, keep its held-out part private, and extend the executor to repository-level tasks with per-task environments.
\item \textbf{Several samples per item.} Record several replies per item, or replies under several batch configurations, to estimate the nondeterministic component of each comparison, and evaluate batch-invariant inference~\cite{he2025nondeterminism} as a way to remove it.
\item \textbf{Complete Experiment C.} Record the bf16, Q8\_0 and ternary arms on HumanEvalFix, commit their reports, and regenerate the tables of Section V-F with \texttt{scripts/\allowbreak{}make\_\allowbreak{}report\_\allowbreak{}tables.\allowbreak{}py}.
\end{enumerate}

\bibliographystyle{IEEEtran}
\bibliography{refs}

\end{document}
